\chapter{Cơ sở lý thuyết và Yêu cầu hệ thống}
\label{ch:background-requirements}

Chương này gồm hai phần bổ trợ cho nhau. Phần thứ nhất trình bày cơ sở lý thuyết và được giới hạn có chủ
đích ở những khái niệm mà thiết kế cùng môi trường kiểm chứng sử dụng trực tiếp: cơ chế tín hiệu của
I\textsuperscript{2}C, chế độ SDR của MIPI I3C Basic, các điều kiện bus, hai kiểu lái Open Drain và
Push-Pull, thủ tục cấp địa chỉ động và Common Command Code. Phần thứ hai chuyển các khái niệm đó thành
yêu cầu hệ thống có thể kiểm chứng; các yêu cầu này là căn cứ để lựa chọn kiến trúc tại
\cref{ch:architecture}, xây dựng môi trường UVM tại \cref{ch:verification} và đánh giá kết quả tại
\cref{ch:results}.

\section{Cơ sở lý thuyết}
\label{sec:bg-theory}

\subsection{Tổng quan bus I\textsuperscript{2}C}
\label{sec:bg-i2c}

I\textsuperscript{2}C là bus nối tiếp đồng bộ hai dây gồm đường dữ liệu SDA và đường xung nhịp SCL. Mọi
thiết bị nối vào bus qua ngõ ra Open Drain: mỗi thiết bị chỉ có khả năng kéo đường tín hiệu xuống mức
thấp, còn mức cao được thiết lập bởi điện trở kéo lên bên ngoài. Cấu trúc wired-AND này cho phép nhiều
thiết bị cùng chia sẻ hai đường tín hiệu và là nền tảng của cơ chế arbitration: khi hai thiết bị phát
đồng thời, thiết bị phát mức thấp chiếm ưu thế, thiết bị còn lại phát hiện giá trị trên bus khác giá trị
mình phát và tự rút lui mà không gây xung đột điện.

Một giao dịch I\textsuperscript{2}C điển hình bắt đầu bằng điều kiện START, tiếp theo là Address Header
gồm địa chỉ 7-bit và bit hướng RnW, sau đó là các byte dữ liệu. Mỗi byte được bên nhận xác nhận bằng một
bit ACK/NACK, và giao dịch kết thúc bằng điều kiện STOP. Trong Fast Mode, tần số SCL tối đa là
400~kHz~\cite{nxp_i2c}.

Hai giới hạn của I\textsuperscript{2}C liên quan trực tiếp đến phạm vi khóa luận. Thứ nhất, vì mức cao
chỉ được tạo bởi điện trở kéo lên, thời gian cạnh lên phụ thuộc vào tích của điện trở kéo lên và điện
dung tải trên bus; đây là nguyên nhân vật lý giới hạn tần số hoạt động. Thứ hai, mỗi thiết bị mang một
Static Address cố định, nên việc tích hợp nhiều thiết bị cùng loại trên một bus đòi hỏi cấu hình chân
địa chỉ hoặc thêm mạch phụ trợ, và hệ thống không có cách chuẩn hóa để phát hiện hay cấp phát địa chỉ
lúc vận hành.

\subsection{MIPI I3C Basic và chế độ SDR}
\label{sec:bg-sdr}

MIPI I3C Basic giữ nguyên hai đường SDA/SCL và khả năng cùng tồn tại với thiết bị
I\textsuperscript{2}C, nhưng định nghĩa lại vai trò thiết bị thành Controller và Target; Controller đang
nắm quyền điều khiển bus được gọi là Active Controller. Trong Single Data Rate Mode (SDR), một bit dữ
liệu được truyền trong mỗi chu kỳ SCL và tần số SCL có thể đạt 12,5~MHz khi hệ thống đáp ứng các tham số
timing tương ứng~\cite{mipi_i3c_basic}.

Cải tiến tốc độ của I3C đến từ việc phân chia lại kiểu lái theo từng pha. Chỉ những pha cần arbitration
hoặc cần phản hồi từ nhiều thiết bị --- Address Header và các bit xác nhận --- còn giữ Open Drain; pha
dữ liệu thông thường được chuyển sang Push-Pull, nhờ đó cạnh tín hiệu không còn phụ thuộc điện trở kéo
lên. Bên cạnh đó, I3C thay bit ACK sau mỗi byte dữ liệu bằng T-Bit với ngữ nghĩa phụ thuộc hướng truyền:
khi Controller ghi, T-Bit mang odd parity của byte vừa phát, cho phép Target phát hiện lỗi đơn bit; khi
Controller đọc, T-Bit là phương tiện để bên phát dữ liệu báo tiếp tục hoặc để phía nhận kết thúc truyền
theo thể thức chuyển giao quyền lái được đặc tả quy định~\cite{mipi_i3c_basic}.

\subsection{START, Repeated START và STOP}
\label{sec:bg-conditions}

Các điều kiện bus của I3C kế thừa từ I\textsuperscript{2}C. START xuất hiện khi SDA chuyển từ cao xuống
thấp trong lúc SCL ở mức cao; STOP là chuyển tiếp ngược lại của SDA khi SCL cao. Repeated START (Sr) có
dạng sóng giống START nhưng xuất hiện khi giao dịch chưa có STOP, cho phép Active Controller mở một pha
địa chỉ mới mà không nhả quyền điều khiển bus. Sau STOP, đặc tả còn phân biệt Bus Free Condition, Bus
Available Condition và Bus Idle Condition theo khoảng thời gian SDA/SCL cùng duy trì mức cao; các mức
phân biệt này chủ yếu phục vụ Hot-Join và IBI. Do hai tính năng đó nằm ngoài phạm vi, thiết kế chỉ cần
tạo và phát hiện đúng START, Sr, STOP và bảo đảm thời gian Bus Free trước giao dịch kế tiếp.

\subsection{Open Drain và Push-Pull}
\label{sec:bg-odpp}

Open Drain biểu diễn mức cao bằng trạng thái high-impedance và nhờ phần tử kéo lên bên ngoài xác lập
điện áp, còn Push-Pull điều khiển chủ động cả hai mức logic. Push-Pull cho cạnh tín hiệu nhanh nhưng
không cho phép nhiều thiết bị cùng phát, vì vậy việc chuyển đổi giữa hai kiểu lái phải theo pha của giao
thức: trong một SDR private transfer, START, Address Header và bit ACK diễn ra ở Open Drain vì Target
cần khả năng phản hồi trên bus; từ sau ACK, các byte dữ liệu cùng T-Bit được truyền ở Push-Pull. Ranh
giới chuyển đổi nằm ngay tại biên của pha xác nhận, do đó phần cứng phải nhận biết chính xác từng pha
thay vì gắn cố định một kiểu lái cho mỗi giao dịch. Đây là một trong những yêu cầu điều khiển trung tâm
mà RTL phải hiện thực.

\subsection{Định dạng SDR private transfer}
\label{sec:bg-frame}

Một SDR private transfer có cấu trúc tuần tự như sau: điều kiện START; Address Header gồm Dynamic
Address 7-bit và bit RnW phát ở Open Drain; bit ACK của Target; các word dữ liệu ở Push-Pull, mỗi byte
đi kèm một T-Bit; cuối cùng là Repeated START nếu còn giao dịch nối tiếp hoặc STOP nếu kết thúc. Với
private write, Active Controller phát từng byte cùng T-Bit odd parity. Với private read, Target phát dữ
liệu và Active Controller dùng T-Bit để duy trì hoặc chấm dứt giao dịch. Từ định dạng này, ba nghĩa vụ
của phần cứng được rút ra trực tiếp: bảo toàn thứ tự bit MSB-first, đếm chính xác số byte đã truyền hoặc
nhận, và tạo STOP hay Sr đúng theo thuộc tính kết thúc ghi trong Command Descriptor.

\subsection{Dynamic Address Assignment và ENTDAA}
\label{sec:bg-daa}

Dynamic Address Assignment (DAA) là thủ tục Active Controller cấp Dynamic Address cho các Target I3C,
thay cho mô hình địa chỉ tĩnh của I\textsuperscript{2}C. Với ENTDAA, Active Controller phát START, gửi
I3C Broadcast Address \texttt{7'h7E} với RnW bằng Write, và sau ACK phát byte CCC \texttt{0x07} để đưa
các Target chưa có địa chỉ vào trạng thái tham gia cấp phát. Mỗi vòng cấp phát bắt đầu bằng Repeated
START và Address Header \texttt{7'h7E} với RnW bằng Read: các Target đủ điều kiện cùng ACK, sau đó phát
đồng thời chuỗi định danh 64-bit gồm PID 48-bit, BCR 8-bit và DCR 8-bit trên Open Drain. Nhờ tính chất
wired-AND, thiết bị có chuỗi định danh nhỏ nhất thắng arbitration mà không cần trọng tài riêng. Active
Controller thu chuỗi 64-bit của thiết bị thắng, phát Dynamic Address 7-bit kèm bit parity, và Target xác
nhận bằng ACK. Chu trình lặp lại cho Target kế tiếp đến khi không còn thiết bị phản hồi Address Header
hoặc đã cấp đủ số lượng yêu cầu, sau đó Active Controller kết thúc bằng STOP~\cite{mipi_i3c_basic}.

\subsection{Common Command Code}
\label{sec:bg-ccc}

Common Command Code (CCC) là cơ chế điều khiển chuẩn hóa giữa Active Controller và Target, được phát
sau Address Header \texttt{7'h7E}. Broadcast CCC áp dụng cho mọi Target trên bus; Direct CCC bổ sung
một pha Repeated START và Address Header của Target cụ thể trước phần dữ liệu lệnh. Tập CCC được hiện
thực trong khóa luận gồm ENEC, DISEC ở cả hai dạng và ENTDAA, như \cref{tab:ccc}; định dạng frame chi
tiết của từng lệnh được liệt kê tại \cref{app:ccc}.

\begin{table}[H]
  \centering
  \caption{Tập CCC được hiện thực.}
  \label{tab:ccc}
  \begin{tabular}{L{3.0cm} C{2.0cm} C{2.6cm} L{5.3cm}}
    \toprule
    \textbf{CCC} & \textbf{Opcode} & \textbf{Dạng} & \textbf{Chức năng trong thiết kế} \\
    \midrule
    ENEC & \texttt{0x00} & Broadcast & Phát event byte đến các Target. \\
    DISEC & \texttt{0x01} & Broadcast & Phát event byte để vô hiệu sự kiện. \\
    ENTDAA & \texttt{0x07} & Broadcast & Khởi động Dynamic Address Assignment. \\
    ENEC & \texttt{0x80} & Direct & Phát event byte đến một Target. \\
    DISEC & \texttt{0x81} & Direct & Phát event byte đến một Target. \\
    \bottomrule
  \end{tabular}
\end{table}

\subsection{So sánh I3C và I\textsuperscript{2}C trong phạm vi thiết kế}
\label{sec:bg-compare}

\Cref{tab:i3c-vs-i2c} tổng hợp các thuộc tính khác biệt có ảnh hưởng trực tiếp đến kiến trúc RTL. I3C
còn định nghĩa nhiều tính năng khác như HDR, IBI hay Hot-Join, nhưng chúng không được dùng làm yêu cầu
cho thiết kế này nên không đưa vào so sánh.

\begin{table}[H]
  \centering
  \caption{Các thuộc tính liên quan của I3C SDR và I\textsuperscript{2}C Fast Mode.}
  \label{tab:i3c-vs-i2c}
  \begin{tabular}{L{4.0cm} L{4.3cm} L{4.6cm}}
    \toprule
    \textbf{Thuộc tính} & \textbf{I\textsuperscript{2}C Fast Mode} & \textbf{I3C SDR} \\
    \midrule
    Đường bus & SDA, SCL & SDA, SCL \\
    Kiểu điều khiển & Open Drain & Open Drain và Push-Pull \\
    Tần số SCL trong phạm vi & Tối đa 400~kHz & Tối đa 12,5~MHz \\
    Địa chỉ & Static Address & Dynamic Address qua DAA \\
    Xác nhận dữ liệu & ACK/NACK sau mỗi byte & T-Bit theo hướng truyền \\
    Điều khiển chuẩn hóa & Không có CCC & Broadcast CCC và Direct CCC \\
    Khả năng cùng tồn tại & Thiết bị I\textsuperscript{2}C & Có thể cùng bus với Target I\textsuperscript{2}C \\
    \bottomrule
  \end{tabular}
\end{table}

Vì cả hai giao thức cùng chạy trên một cặp SDA/SCL, thiết kế phải tạo được hai bộ timing khác nhau từ
cùng một clock hệ thống. \Cref{tab:timing} liệt kê các tham số chính dùng để xác lập cấu hình mặc định;
giá trị I3C lấy từ MIPI I3C Basic v1.1.1, giá trị I\textsuperscript{2}C lấy từ UM10204
Rev.~7.0~\cite{mipi_i3c_basic,nxp_i2c}. Trong RTL, mỗi tham số được biểu diễn bằng số chu kỳ clock hệ
thống và có thể cấu hình qua CSR.

\begin{table}[H]
  \centering
  \caption{Các tham số timing chính dùng để xác lập cấu hình mặc định.}
  \label{tab:timing}
  \begin{tabular}{L{4.4cm} C{3.0cm} C{3.0cm} L{2.5cm}}
    \toprule
    \textbf{Tham số} & \textbf{I3C SDR} & \textbf{I\textsuperscript{2}C FM} & \textbf{Nguồn} \\
    \midrule
    Tần số SCL tối đa & 12,5~MHz & 400~kHz & Đặc tả bus \\
    Thời gian SCL thấp tối thiểu & 24~ns & 1,3~$\mu$s & Đặc tả bus \\
    Thời gian SCL cao tối thiểu & 24~ns & 0,6~$\mu$s & Đặc tả bus \\
    Data setup tối thiểu & 3~ns (PP) & 100~ns & Đặc tả bus \\
    \bottomrule
  \end{tabular}
\end{table}

\subsection{Thuật ngữ UVM 1.2}
\label{sec:bg-uvm}

UVM là thư viện lớp SystemVerilog dùng để tổ chức Testbench có khả năng tái sử dụng. Agent đóng gói
Sequencer, Driver và Monitor cho một interface: Sequence tạo transaction, Driver chuyển transaction
thành tín hiệu, còn Monitor thực hiện chiều ngược lại một cách độc lập với Driver. Scoreboard kiểm tra
tính đúng đắn ở mức transaction bằng cách đối chiếu dữ liệu quan sát với reference model. Functional
Coverage đo các tình huống đã thực sự xảy ra trong simulation, còn SVA (SystemVerilog Assertions) kiểm
tra thuộc tính theo chu kỳ ở mức tín hiệu. Virtual Sequence điều phối nhiều Agent thông qua Virtual
Sequencer, nhờ đó một kịch bản có thể đồng thời cấu hình DUT qua bus thanh ghi và điều khiển phản hồi
trên bus I3C~\cite{uvm12_ug,uvm12_ref}. Chương này chỉ xác lập thuật ngữ; topology cụ thể của môi
trường được trình bày tại \cref{ch:verification}.

\section{Yêu cầu hệ thống}
\label{sec:req}

\subsection{Yêu cầu chức năng}
\label{sec:req-func}

Từ cơ sở lý thuyết trên và phạm vi đã xác định ở \cref{ch:introduction}, yêu cầu chức năng của thiết kế được
phát biểu thành mười mục F1--F10 trong \cref{tab:req-func}. Mỗi mục tương ứng một nhóm hành vi có thể
kiểm chứng độc lập và được dùng làm đơn vị ánh xạ sang cơ chế kiểm chứng ở \cref{ch:verification}.

\begin{table}[H]
  \centering
  \caption{Yêu cầu chức năng của I3C Active Controller.}
  \label{tab:req-func}
  \begin{tabular}{C{1.4cm} L{11.8cm}}
    \toprule
    \textbf{Mã} & \textbf{Yêu cầu} \\
    \midrule
    F1 & Thực hiện SDR private write đến Target I3C, gồm Address Header, ACK, payload và T-Bit. \\
    F2 & Thực hiện SDR private read từ Target I3C và kết thúc đúng theo độ dài hoặc T-Bit. \\
    F3 & Hỗ trợ Immediate Data Transfer với tối đa bốn byte trong Command Descriptor. \\
    F4 & Hỗ trợ Broadcast CCC và Direct CCC cho ENEC/DISEC. \\
    F5 & Thực hiện ENTDAA, lưu PID/BCR/DCR và cấp Dynamic Address theo các mục DAT được chỉ định. \\
    F6 & Thực hiện private read và write với Target I\textsuperscript{2}C Fast Mode. \\
    F7 & Cung cấp Command Queue, Response Queue, TX/RX Data Buffer và DAT cho luồng phần mềm--phần cứng. \\
    F8 & Trả Response Descriptor chứa transaction ID, số byte và error status. \\
    F9 & Hỗ trợ hardware abort, software reset, backpressure và phát hiện tràn/thiếu dữ liệu. \\
    F10 & Cho phép cấu hình timing I3C/I\textsuperscript{2}C và trạng thái enable qua CSR. \\
    \bottomrule
  \end{tabular}
\end{table}

\subsection{Các tính năng ngoài phạm vi}
\label{sec:req-oos}

Ranh giới phạm vi được phát biểu tường minh trong \cref{tab:req-oos} nhằm hai mục đích: tránh suy diễn
yêu cầu ngầm trong quá trình thiết kế, và làm chuẩn khi xây dựng coverage --- một tính năng ngoài phạm
vi không được tính vào coverage goal, cũng không bị đánh dấu là tình huống cấm chỉ vì RTL không hỗ trợ.

\begin{table}[H]
  \centering
  \caption{Các tính năng không thuộc phạm vi hiện thực.}
  \label{tab:req-oos}
  \begin{tabular}{L{4.0cm} L{9.1cm}}
    \toprule
    \textbf{Tính năng} & \textbf{Ranh giới phạm vi} \\
    \midrule
    HDR & Không hiện thực HDR-DDR, HDR-TSP hoặc HDR-TSL. \\
    In-Band Interrupt & Không tiếp nhận hoặc điều phối IBI. \\
    Hot-Join & Không xử lý yêu cầu Target tham gia bus khi hệ thống đang hoạt động. \\
    Secondary Controller & Không chuyển giao vai trò Active Controller. \\
    Target mode & IP chỉ hoạt động ở vai trò Active Controller. \\
    Full HCI compliance & Chỉ giữ các Descriptor, Queue và CSR cần cho phạm vi. \\
    FPGA validation & Không có dữ liệu synthesis, timing closure hoặc kiểm thử trên board. \\
    \bottomrule
  \end{tabular}
\end{table}

\subsection{Chỉ tiêu hiệu năng và cấu hình}
\label{sec:req-perf}

\begin{table}[H]
  \centering
  \caption{Chỉ tiêu hiệu năng và cấu hình của hệ thống.}
  \label{tab:req-perf}
  \begin{tabular}{C{1.4cm} L{6.8cm} L{5.0cm}}
    \toprule
    \textbf{Mã} & \textbf{Chỉ tiêu} & \textbf{Giá trị mục tiêu} \\
    \midrule
    P1 & Clock hệ thống & Tối thiểu 333~MHz \\
    P2 & I3C SDR SCL & Có thể cấu hình đến 12,5~MHz \\
    P3 & I\textsuperscript{2}C Fast Mode SCL & Có thể cấu hình đến 400~kHz \\
    P4 & Độ sâu DAT & 32 mục \\
    P5 & Độ rộng dữ liệu giao diện thanh ghi & 32 bit \\
    \bottomrule
  \end{tabular}
\end{table}

Chỉ tiêu P1 xuất phát từ độ phân giải timing: để biểu diễn các tham số cỡ 24~ns của I3C bằng số nguyên
chu kỳ với sai số chấp nhận được, chu kỳ clock hệ thống cần nhỏ hơn nhiều lần tham số nhỏ nhất; thiết
kế chọn 3~ns làm chu kỳ danh định. Các chỉ tiêu còn lại cố định kích thước những cấu trúc mà phần mềm
nhìn thấy, làm cơ sở cho bản đồ thanh ghi tại \cref{app:csr}.

\subsection{Giao diện phần mềm}
\label{sec:req-sw}

Software-visible register interface sử dụng địa chỉ 12-bit, dữ liệu 32-bit, hai tín hiệu yêu cầu
đọc/ghi và một tín hiệu hoàn tất. Phần mềm cấu hình enable, abort, software reset và timing qua CSR;
ghi hai DWORD liên tiếp vào cổng Command Queue để tạo Command Descriptor 64-bit; nạp payload vào TX
Data Buffer; và sau khi giao dịch kết thúc, đọc RX Data Buffer cùng Response Queue. DAT lưu Static
Address hoặc Dynamic Address kèm thuộc tính loại Target cho từng thiết bị. Giao diện không hỗ trợ burst
hoặc transaction ID ở mức bus thanh ghi. Bản đồ thanh ghi đầy đủ được đặt tại \cref{app:csr}, định dạng
Descriptor tại \cref{app:desc}.

\subsection{Yêu cầu kiểm chứng}
\label{sec:req-verif}

Môi trường kiểm chứng phải tạo được stimulus cho toàn bộ yêu cầu F1--F10, bao gồm nominal case và các
trường hợp ACK/NACK, độ dài biên, FIFO full/empty, reset và abort. Scoreboard phải kiểm tra payload,
Address Header, T-Bit, Response Descriptor và side effect của CSR. SVA phải kiểm tra handshake, FSM,
reset, Queue và các điều kiện bus quan trọng. Functional Coverage phải ghi nhận loại giao dịch, hướng
truyền, địa chỉ, độ dài, response/error, trạng thái Buffer và các corner case. Một regression log chỉ
được xem là đạt khi phần tổng kết UVM có \texttt{UVM\_ERROR=0} và \texttt{UVM\_FATAL=0}; coverage được
đánh giá riêng và không được dùng thay cho kết quả kiểm tra đúng/sai.
